%B1===============================[=.............................]===============================%B1
\def\BeightEN{}
\def\BoneLT{$\dfrac{20\cdot 24}{2\cdot 0 + 2\cdot 4}=$}
\def\BeightPL{}
\def\BeightRU{}

\def\rBone{D}
\def\aBone{\answer{12}{30}{48}{60}{120}}

%B2===============================[==............................]===============================%B2
\def\BtwoEN{Dina has set up her three bricks on the floor behind a wall. When seen from the front, the bricks look like this .
How do the bricks look from the back ?}
\def\BtwoLT{Vytis pastatė tris bokštelius už dėžutės (žr.~pav.~dešinėje). Kaip bokšteliai atrodo iš priešingos pusės?}

\def\rBtwo{A}
\def\aBtwo{\gransweranswer[0.8]{B01-3}{B01-2}{B01-4}{B01-5}{B01-6}}
\def\aBtwoSP{\granswer[0.8]{B01-3}{B01-2}{B01-4}{B01-5}{B01-6}}


%B3===============================[===...........................]===============================%B3
\def\BthreeEN{The picture shows the first few squares of a hopping game. Every fourth square in the game has the same
image in it. Mia is playing the game. In which of the following squares will Mia land only on her right foot?}
\def\BthreeLT{Milda šokinėja ant šaligatvio plytelių, kaip parodyta paveikslėlyje:
ant pirmosios plytelės šoka abiem kojom, ant antrosios -- tik kaire, ir t.~t. vis kartodama pirmus keturis šuolius.
%Milda šokinėja ant šaligatvio plytelių: ant pirmos plytelės šoka abiem kojom, ant antrosios -- tik kaire, ant trečiosios -- vėl abiem, ant ketvirtosios -- tik dešine,
%ir taip tęsia toliau.
%o toliau šiuos šuolius vis kartoja: į penktąjį langelį vėl įšoka kaire koja, į šeštąjį -- vėl abiem, ir t.~t.
Ant kelintos plytelės Milda užšoks tik dešine koja?
\vspace{-5pt}\begin{center}
                  \includegraphics[scale=0.8]{B02-1} \vspace{-5pt}
\end{center}
}
\def\rBthree{C}
\def\aBthreeEN{\answer{the 10th}{the 15th}{the 20th}{the 22nd}{the 23rd}}
\def\aBthreeLT{\answer{Ant 10-os}{Ant 15-os}{Ant 20-os}{Ant 22-os}{Ant 23-ios}}

%B4===============================[====..........................]===============================%B4

\def\BfourEN{Sasha created a secret alphabet. He writes ``basil'' as
\includegraphics[scale=0.5]{B03-3}  and ``red'' as \includegraphics[scale=0.6]{B03-4} . How does he write ``bread''?}
\def\BfourLT{Ona sukūrė slaptą abėcėlę, kurioje kiekvieną lietuviškos abėcėlės raidę atitinka slaptas simbolis.
Savo katės vardą PILKĖ ji užrašo taip: \includegraphics[scale=0.35]{B03-3},
    o vardą ONA -- taip: \includegraphics[scale=0.35]{B03-4}. Kaip ji užrašo žodį PONIA?}
\def\BfourEXP{Ona sukūrė slaptą abėcėlę, kurioje kiekvieną lietuviškos abėcėlės raidę atitinka slaptas simbolis.
Savo katės vardą PILKĖ ji užrašo taip: \includegraphics[scale=0.5]{B03-3},
    o vardą ONA -- taip: \includegraphics[scale=0.5]{B03-4}. Kaip ji užrašo žodį PONIA?}
\def\rBfour{C}
\def\aBfour{\granswer[0.35]{B03-5}{B03-6}{B03-7}{B03-8}{B03-9}}
\def\aBfourEXP{\granswer[0.5]{B03-5}{B03-6}{B03-7}{B03-8}{B03-9}}



%B5===============================[=====.........................]===============================%B5

\def\BfiveEN{Which of the strips should be placed in the space in the picture so that each child is connected to a different kite?}
\def\BfiveLT{Abu vaikai paveikslėlyje laiko po aitvarą. Kuri iš iškirptų juostelių yra paveikslėlio dalis?}
%arba: dalis uždengta juostele. kuri? ir pan.}

\def\rBfive{D}
\def\aBfive{\granswerbenjaminfive[0.7]{B04-1}{B04-2}{B04-3}{B04-4}{B04-5}}
\def\aBfiveSP{\granswer[0.7]{B04-1}{B04-2}{B04-3}{B04-4}{B04-5}}

%B6===============================[======........................]===============================%B6
\def\BsevenEN{Alina folds the image below along the dashed line. Which of the following squares folds onto an identical one?}
\def\BsevenLT{Paveikslėlis sudarytas iš skirtingų rūšių langelių.
Rasa perlenkė paveikslėlį per punktyrinę liniją.
Kurios rūšies du langeliai perlenkus sutapo?
}

\def\rBseven{A}
\def\aBseven{\granswer[0.6]{B05-3}{B05-2}{B05-4}{B05-5}{B05-6}}

%B7===============================[=======.......................]===============================%B7

\def\BeightEN{Mona wants to draw the figure shown without lifting up her pencil from the paper.
The lengths of the three segments are given. What is the shortest total length she could draw?}
\def\BeightLT{Neatitraukdamas rankos, Tomas pieštuku nubrėžė figūrą, kurią sudaro trys atkarpos.
Paveikslėlyje parodyta ši figūra ir visų atkarpų ilgiai.
Kokį trumpiausią kelią galėjo popieriumi nueiti pieštukas, Tomui brėžiant figūrą?}

\def\rBeight{B}
\def\aBeight{\answer{$6\ \textup{cm}$}{$7\ \textup{cm}$}{$8\ \textup{cm}$}{$9\ \textup{cm}$}{$10\ \textup{cm}$}}



%B8===============================[========......................]===============================%B8
%\def\BeightEN{There are six boxes on a truck as shown.
%A worker puts them on the floor. He takes one box at a time, provided that box does not have another box on top of it.
%He places his box on the ground or on top of another box. Which of the following stacks could he not build?}
%\def\BeightLT{Sunkvežimyje sudėtos šešios dėžės (žr. pav. šalia).
%Darbininkas ima po vieną dėžę, ant kurios nėra kitos dėžės, ir deda ją į naują krūvą.
%Kaip tikrai negali atrodyti iškrautos dėžės?}
%\def\BeightPL{}
%\def\BeightRU{}
%
%\def\rBeight{C}
%\def\aBeight{\granswer[0.5]{B08-2}{B08-3}{B08-4}{B08-5}{B08-6}}

\def\BsixEN{}
\def\BsixLT{Vieną rytą Ūla nusipirko 14 sausainių. Kiekvieną popietę ji suvalgydavo po pusę turimų sausainių.
Kiekvieną vakarą Ūlos šuo suėsdavo po vieną sausainį. Galiausiai sausainių nebeliko. Kiek iš viso sausainių suėdė Ūlos šuo?}
\def\BsixPL{224 minuty temu była godzina 20 : 24. Która godzina jest teraz?}


\def\rBsix{E}
\def\aBsix{\answer{7}{6}{5}{4}{3}}



%B9===============================[=========.....................]===============================%B9
\def\BnineEN{Pieter has a package of $445 \mbox{g}$ and the following eight weights: \includegraphics{B09-1}.
He put the package on the scale, as shown. What is the minimum number of weights he needs to balance the scale?}
\def\BnineLT{Petras turi aštuonis svarelius, kurių masė nurodyta gramais: \includegraphics[scale=0.8]{B09-1}.
Jis padėjo $445\ \mbox{g}$ sveriantį paketą ant lėkštinių svarstyklių.
Kiek mažiausiai svarelių Petrui reikia padėti ant svarstyklių, kad jos taptų pusiausviros?
}
\def\BnineEXP{Petras turi aštuonis svarelius, kurių masė nurodyta gramais:

%\begin{center}
\hspace{20pt}    \includegraphics[scale=1]{B09-1}.
%\end{center}

Jis padėjo $445\ \mbox{g}$ sveriantį paketą ant lėkštinių svarstyklių.
Kiek mažiausiai svarelių Petrui reikia padėti ant svarstyklių, kad jos taptų pusiausviros?
}
\def\BninePL{}
\def\BnineRU{}

\def\rBnine{B}
\def\aBnine{\answer{\(2\)}{\(3\)}{\(4\)}{\(5\)}{\(6\)}}

%B10===============================[==========....................]===============================%B10
\def\BtenEN{The rooms in the hotel are numbered in ascending order, starting from 1.
No number is omitted. Kangaroo counted the digits in the rooms and found digit 2 14 times and digit 5 3 times.
What is the largest number of rooms there can be in the hotel?}
\def\BtenLT{Viešbutyje visi kambariai sunumeruoti iš eilės pradedant nuo 1. Tauras suskaičiavo kambarių numerių skaitmenis ir gavo,
kad skaitmuo 2 pasikartoja 14 kartų, o skaitmuo~5 pasikartoja 3 kartus. Kiek daugiausiai kambarių gali būti viešbutyje?}
\def\BtenPL{}
\def\BtenRU{}

\def\rBten{C}
\def\aBten{\answer{25}{26}{34}{35}{41}}

%B11===============================[===========...................]===============================%B11
\def\BuoneEN{Two identical rectangles, each with an area of 18, overlap to form a new rectangle, as shown.
The new rectangle can be divided into three identical squares.
What is the area of the new rectangle?}
\def\BuoneLT{Ant stalo guli 24 figūrėlės. Kiekviena iš jų yra stačiakampis arba rombas. Stačiakampių yra 15, o rombų
irgi yra 15. Kiek figūrėlių yra kvadratai?}
\def\BuonePL{}
\def\BuoneRU{}

\def\rBuone{B}
\def\aBuone{\answer{0}{6}{12}{15}{18}}

%B12===============================[============..................]===============================%B12
\def\ButwoEN{Która z poniższych liczb całkowitych jest o dwa mniejsza od wielokrotności liczby 10, o dwa
większa od kwadratu liczby naturalnej i dwa razy większa od pewnej liczby pierwszej?}
\def\ButwoLT{Kuris iš atsakymuose išvardytų skaičių yra dviem mažesnis už 10 kartotinį, dviem didesnis už natūraliojo skaičiaus kvadratą ir dvigubai didesnis už pirminį skaičių?}
\def\ButwoPL{Która z poniższych liczb całkowitych jest o dwa mniejsza od wielokrotności liczby 10, o dwa
większa od kwadratu liczby naturalnej i dwa razy większa od pewnej liczby pierwszej?}
\def\ButwoRU{}

\def\rButwo{D}
\def\aButwo{\answer{98}{78}{58}{38}{18}}

%B13===============================[=============.................]===============================%B13
\def\ButhreeEN{Rosa draws several identical rectangles to make the following picture. The width and the height of the picture are $45\ \text{cm}$ and $30\ \text{cm}$ respectively. What is the area of one rectangle?}
\def\ButhreeLT{
%Saulė iš vienodų kaladėlių pastatė 
Saulė iš keleto lygių stačiakampių sudėliojo figūrą (žr. pav.).
    Figūros pagrindo ilgis yra $45\ \text{cm}$, o visos figūros aukštis yra $30\ \text{cm}$. Koks yra vieno stačiakampio plotas?}
\def\ButhreePL{}
\def\ButhreeRU{}

\def\rButhree{E}
\def\aButhree{\answer{$24\ \text{cm}^2$}{$27\ \text{cm}^2$}{$30\ \text{cm}^2$}{$33\ \text{cm}^2$}{$36\ \text{cm}^2$}}

%B14===============================[==============................]===============================%B14
\def\BufourEN{Each of the 16 circles shown contains a number. Numbers in neighbouring circles differ by 1.
One of the circles contains the number 5 and another one contains 13. How many different numbers are written in the 16 circles?}
\def\BufourLT{Į kiekvieną iš 8 skritulių įrašyta po skaičių.
Bet kurie du gretimuose skrituliuose esantys skaičiai skiriasi vienetu.
%Gretimuose apskritimuose skaičiai skiriasi vienetu.
Į vieną iš skritulių įrašytas skaičius 5, o į kitą -- skaičius 9.
Kiek skirtingų skaičių įrašyta?}
\def\BufourPL{}
\def\BufourRU{}

\def\rBufour{B}
\def\aBufour{\answer{4}{5}{6}{7}{8}}

%B15===============================[===============...............]===============================%B15
\def\BufiveEN{The diagram shows two large squares with the same area. Part of each square is shaded, as shown. In the first square, the midpoints of adjacent sides are joined. In the second square, four smaller squares all with side-lengths equal to a third of the side-lenght of the large squrare are shaded. The area shaded in the first square is $9$. What is the area shaded in the second square?}
\def\BufiveLT{Diagramoje pavaizduoti du vienodo ploto kvadratai $A$ ir $B$.
Kvadrato $A$  gretimų kraštinių vidurio taškai yra sujungti atkarpomis.
Kvadrato $B$  kampuose pažymėti 4 kvadratėliai, kurių kraštinių ilgiai lygūs vienai trečiajai didžiojo kvadrato kraštinės ilgio.
%$B$ kvadrato viduje esančių 4 mažesnių nuspalvintų kvadratų kraštinių ilgiai lygūs vienai trečiajai didžiojo kvadrato kraštinės ilgio.
Kvadrato $A$ pilkosios dalies plotas lygus $9$. Kam lygus kvadrato  $B$ pilkosios  dalies plotas?}
\def\BufivePL{}
\def\BufiveRU{}

\def\rBufive{B}
\def\aBufive{\answer{4}{8}{9}{10}{12}}

%B16===============================[================..............]===============================%B16
\def\BusixEN{The Braille system for blind people, when written down, has the digits 0 to 9 represented by a set of black or white dots, as shown. \newline \centerline{} \newline How many different two-digit numbers contain exactly five black dots?}
\def\BusixLT{Skaitmenys nuo 0 iki 9, užrašyti Brailio raštu, paveikslėlyje pavaizduoti kaip juodų ir baltų taškų kombinacijos.
\vspace{-5pt}
\begin{center} \includegraphics[scale=0.3]{B16-1} \end{center}
\vspace{-8pt}
Kiek dviženklių skaičių turi lygiai 5 juodus taškus?
}
\def\BusixPL{}
\def\BusixRU{}

\def\rBusix{C}
\def\aBusix{\answer{16}{18}{30}{32}{34}}

%B17===============================[=================.............]===============================%B17
\def\BusevenEN{The figure below shows a beehive with 16 cells. Some of the cells contain honey. The number in each cell indicates how many of its neighbouring cells contain honey. Two cells are neighbours if they share a common edge. How many cells in the beehive contain honey?}
\def\BusevenLT{Paveikslėlyje pavaizduotas korys su 16 akučių. Kai kurios iš akučių užpildytos medumi. Akutės yra kaimynės, jei turi lygiai vieną bendrą kraštinę.
Skaičius akutėje lygus medumi užpildytų jos kaimynių skaičiui. Kiek korio akučių yra užpildytos medumi?}
\def\BusevenPL{}
\def\BusevenRU{}

\def\rBuseven{C}
\def\aBuseven{\answer{7}{8}{9}{10}{11}}

%B18===============================[==================............]===============================%B18
\def\BueightEN{Annie wants to place the numbers 1 to 10 in the circles in the diagram with one number in each circle. She wants the sum of the numbers in any four circles that are in a straight line, for example the four grey ones, to be 23. What number must she place in the circle containing the question mark?}
\def\BueightLT{
Rusnė ketina visus skaičius nuo 1 iki 10 surašyti į skirtingus skritulius  (žr.~pav.)
%    Rusnė nori į kiekvieną skritulį paveikslėlyje įrašyti po skaičių nuo 1 iki 10.
Ji siekia, kad bet kurių keturių vienoje tiesėje esančių skritulių skaičių suma būtų 23.
Kokį skaičių ji turi įrašyti į skritulį su klaustuku?}
\def\BueightPL{Zuzia chce umieścić w kołach na rysunku obok liczby od 1 do 10, po jednej w każdym kole. Chce, aby suma liczb w dowolnych czterech kołach ułożonych wzdłuż linii prostej wynosiła 23. Jaką liczbę powinna umieścić w kole ze znakiem zapytania?}
\def\BueightRU{}

\def\rBueight{D}
\def\aBueight{\answer{4}{5}{6}{7}{8}}

%B19===============================[===================...........]===============================%B19
\def\BunineEN{Christian has cut four small squares from the corners of the larger square, so that the remaining area is
half of the area of the original square. The side-lenghts of the small squares are shown in the diagram.
What is the perimeter of the remaining shape?}
\def\BunineLT{Barbora nuo kvadratino popieriaus lapo kampų nukirpo keturis kvadratėlius. Kvadratėlių kraštinių ilgiai
centimetrais nurodyti paveikslėlyje.
Likusios figūros plotas yra lygus pusei pradinio kvadrato ploto. Koks yra šios figūros perimetras?}
\def\BuninePL{}
\def\BunineRU{}

\def\rBunine{B}
\def\aBunine{\answer{$36\ \textup{cm}$}{$40\ \textup{cm}$}{$44\ \textup{cm}$}{$48\ \textup{cm}$}{$52\ \textup{cm}$}}

%B20===============================[====================..........]===============================%B20
\def\ButenEN{Ria wants to complete the puzzle shown so that each row and each column contain the numbers 1, 2, 3 and 4 exactly once.
She wants to place the numbers so that the greater than and less symbols ($>$ and $<$) give a correct relationship
between the two values either side of them. The symbols work in all directions, as shown in the example:
\newline \centerline{}\newline What number should she place in the gray circle?}
\def\ButenLT{
Vaiva ketina į kiekvieną skritulį (žr. pav.) įrašyti po skaičių taip,
%Rasa nori išspręsti paveikslėlyje parodytą galvosūkį taip,
kad kiekviename stulpelyje ir kiekvienoje eilutėje
skaičiai 1, 2, 3, 4 pasikartotų lygiai po vieną kartą ir būtų teisingos penkios nurodytos nelygybės ($>$, $<$, $\wedge$ ir $\vee$).
%Nelygybės veikia visomis kryptimis,
%kaip parodyta pavyzdyje:
%\begin{center} \includegraphics[scale=1]{B20-1} \end{center}
Kokį skaičių Vaiva gali įrašyti į pilkąjį skritulį?}
\def\ButenPL{}
\def\ButenRU{}

\def\rButen{A}
\def\aButenEN{\answer{1}{2}{3}{4}{2 or 3}}
\def\aButenLT{\answer{1}{2}{3}{4}{2 arba 3}}
\def\aButenPL{\answer{}{}{}{}{}}
\def\aButenRU{\answer{}{}{}{}{}}

%B21===============================[=====================.........]===============================%B21
\def\BdoneEN{There are three identical special dice on the table. What is the sum of the numbers on the faces that touch the table?}
\def\BdoneLT{Ant stalo guli trys vienodi kubeliai. Kokia yra trijų apatinių sienų skaičių suma?}
\def\BdonePL{}
\def\BdoneRU{}

\def\rBdone{C}
\def\aBdone{\answer{26}{40}{43}{47}{56}}

%B22===============================[======================........]===============================%B22
\def\BdtwoEN{The diagram shows four touching rectangles. What is the area of the shaded rectangle? \newline \centerline{}}
\def\BdtwoLT{Koks yra paveikslėlyje pavaizduoto pilkojo stačiakampio plotas?}
\def\BdtwoPL{}
\def\BdtwoRU{}

\def\rBdtwo{E}
\def\aBdtwoEN{\answer{$12\ \textup{cm}^2$}{$14\ \textup{cm}^2$}{$16\ \textup{cm}^2$}{$18\ \textup{cm}^2$}{$20\ \textup{cm}^2$}}
\def\aBdtwoLT{\answer{$12\ \textup{cm}^2$}{$14\ \textup{cm}^2$}{$16\ \textup{cm}^2$\\}{$18\ \textup{cm}^2$}{$20\ \textup{cm}^2$}}
\def\aBdtwoPL{\answer{}{}{}{}{}}
\def\aBdtwoRU{\answer{}{}{}{}{}}

%B23===============================[=======================.......]===============================%B23
\def\BdfourEN{The figure shows the plan of the seven train routes of a small town.
The circles indicate the stations. Martin wants to paint the lines in such a way that if two lines share a common station,
    then they are painted with different colours. What is the smallest number of colours that he can use?}
\def\BdfourLT{
Dešinėje pavaizduota septynių Kengūrijos traukinių maršrutų schema. Skrituliai žymi stotis.
Per stotis traukiniai visada važiuoja tiesiai.
%Motiejus nori nuspalvinti kiekvieną traukinio liniją taip, kad jei dvi linijos turi bendrą stotį, jos nuspalvintos skirtingomis spalvomis.
Motiejus ketina bendrą stotį turinčių traukinių maršrutus nuspalvinti skirtingomis spalvomis.
    Kiek mažiausiai spalvų panaudos Motiejus?
}
\def\BdfourPL{}
\def\BdfourRU{}

\def\rBdfour{A}
\def\aBdfour{\answer{3}{4}{5}{6}{7}}

%B24===============================[========================......]===============================%B24
\def\BdthreeEN{Dimitri wants to fold the net shown to make a cube.
He wants the triangles that touch the edges of neighbouring faces of the cube to be shaded the same.
How should he shade the triangles of the unshaded square in the net?}
\def\BdthreeLT{Vėjas nori iš dešinėje pavaizuotos išklotinės išlankstyti kubą.
Kaip jis nuspalvins baltąją sieną, jei gretimų sienų bendrą kraštinę turintys trikampėliai turi būti tos pačios
spalvos?
%Kaip turėtų atrodyti baltoji siena,
%    jei gretimų sienų bendrą kraštinę turintys
%trikampėliai yra tos pačios spalvos?
}
\def\BdtreePL{}
\def\BdthreeRU{}

\def\rBdthree{B}
\def\aBdthree{\granswer[0.5]{B24-8}{B24-9}{B24-10}{B24-11}{B24-12}}

%B25===============================[=========================.....]===============================%B25
\def\BdfiveEN{Simon takes four cups out of the cupboard and puts them randomly on the four saucers.
\includegraphics{B25-1}
Which statement is correct?}
\def\BdfiveLT{Ugnius paėmė keturis puodelius su keturiomis lėštutėmis ir atsitiktine tvarka uždėjo puodelius ant lėkštučių.
%r atsi ir padėjo juos atsitiktinai ant keturių jų lėkštučių.
\vspace{-13pt}
\begin{center}
    \includegraphics[scale=0.6]{B25-1}
\end{center}
\vspace{-5pt}
Kuris teiginys teisingas?}
\def\BdfivePL{}
\def\BdfiveRU{}

\def\rBdfive{D}
\def\aBdfiveEN{\answer{It is certain that none of the 4 cups stands on its matching saucer.}
{It is certain that exactly 1 cup stands on its matching saucer.}
{It is impossible for exactly 2 cups to stand on its matching saucer.}
{It is impossible for exactly 3 cups to stand on its matching saucer.}
{It is impossible for all 4 cups to stand on its matching saucer.}}
\def\aBdfiveLT{\answer{Tikrai nė vienas iš 4 puodelių nestovi ant savo lėkštutės\\}
{Lygiai 1 puodelis tikrai stovi ant savo lėkštutės\\}
{Neįmanoma, kad lygiai 2 puodeliai stovėtų ant savo lėkštučių\\}
{Neįmanoma, kad lygiai 3 puodeliai stovėtų ant savo lėkštučių\\}
{Neįmanoma, kad visi 4 puodeliai stovėtų ant savo lėkštučių}}
\def\aBdfivePL{\answer{}{}{}{}{}}
\def\aBdfiveRU{\answer{}{}{}{}{}}

%B26===============================[==========================....]===============================%B26
\def\BdsixEN{A cube with the filled in numbers is given.
Mary wants to write the numbers 1 to 8 on the vertices of the cube.
She wants the sum of the numbers of the vertices around each face to be the same.
She has already written the numbers 6, 7 and 8, as shown. What number should she write on the vertex marked with the questionmark?}
\def\BdsixLT{Joris ketina visose kubo viršūnėse įrašyti skirtingus skaičius nuo 1~iki 8 taip, kad kiekvienos sienos viršūnių skaičių suma būtų ta pati.
Jis jau įrašė skaičius 6,
7 ir 8. Kokį skaičių jis turi įrašyti į klaustuku pažymėtą viršūnę?}
\def\BdsixPL{}
\def\BdsixRU{}

\def\rBdsix{C}
\def\aBdsix{\answer{1}{2}{3}{4}{5}}

%B27===============================[===========================...]===============================%B27
\def\BdsevenEN{A grandmother has some candies. She decides to divide them up amongst her grandchildren so
that each has a bag containing same number of candies. She puts the largest possible number of
candies in each bag and, when she is done, she sees that there are 20 candies in each bag and 12
candies are left over. What is the smallest possible number of candies she could have?}
\def\BdsevenLT{Močiutė turi saldainių ir nori, kad visi anūkai jų gautų po lygiai.
Ji apskaičiavo, kad kievienam anūkui gali duoti daugiausiai po 20 saldainių, ir dar lieka 12 saldainių.
Kiek mažiausiai saldainių ji turi?}
\def\BdsevenPL{}
\def\BdsevenRU{}

\def\rBdseven{C}
\def\aBdseven{\answer{52}{232}{272}{411}{432}}

%B28===============================[============================..]===============================%B28
\def\BdeightEN{Dan plans to cut a rope into 12 equal pieces and marks points where he needs to cut.
Muhammad plans to cut the same rope into 16 equal pieces and marks points where he needs to cut.
Then Maya cuts the rope at all the marked points. How many pieces does Maya get?}
\def\BdeightLT{Jonas planuoja sukarpyti virvę į 12 lygių dalių ir pažymėjo kirpimo vietas. Kostas
planuoja sukarpyti tą pačią virvę į 16 lygių dalių ir taip pat pažymėjo kirpimo vietas.
Jų jaunesnioji sesuo Mėja rado virvę ir sukarpė ją per visas pažymėtas vietas. Kiek virvagalių gavo Mėja?
}
\def\BdeightPL{}
\def\BdeightRU{}

\def\rBdeight{A}
\def\aBdeight{\answer{24}{25}{27}{28}{29}}

%B29===============================[=============================.]===============================%B29
\def\BdnineEN{Emma is playing with the seven caterpillar puzzle pieces shown.
She wants to build a caterpillar that has one head, one tail and either one, two or three puzzle pieces in between.
How many different caterpillars could Emma build?}
\def\BdnineLT{Ieva žaidžia su septyniomis dėlionės detalėmis. Dėlionėje yra dvi galvos, dvi uodegos ir
trys vidurinės detalės.
\vspace{-5pt}
\begin{center}
    \includegraphics[scale=0.6]{B29-1}
\end{center}
\vspace{-5pt}
Ji nori sudaryti vikšrą, kuris turėtų vieną galvos detalę,
vieną uodegos detalę ir vieną, dvi arba tris detales per vidurį.
    Kiek yra skirtingų vikšrų, kuriuos Ieva gali sudaryti?
%    Keliais būdais Ieva gali sukonstruoti vikšrą?
}
\def\BdninePL{}
\def\BdnineRU{}

\def\rBdnine{E}
\def\aBdnine{\answer{10}{14}{16}{18}{20}}

%B30===============================[==============================]===============================%B30
\def\BdtenEN{Ava writes a three-digit number on the whiteboard. Then Brandon writes a fourth digit to the right of
the previous ones. He says ``Look! The number increased by 2024''. What digit did Brandon write?}
\def\BdtenLT{Rūta lentoje užrašė triženklį skaičių. Mėta prirašė ketvirtą skaitmenį iš dešinės ir pasakė: „Skaičius padidėjo
2024 vienetais.“ Kokį skaitmenį prirašė Mėta?}
\def\BdtenPL{}
\def\BdtenRU{}

\def\rBdten{D}
\def\aBdten{\answer{2}{3}{4}{8}{9}}

